% [cc] v22 reference guide, forked from the v21 guide; v22 drafts were reverted, so it matches v21.
% Source of truth: corne_prod_v22.keymap and corne_prod_v22_notation.md.
\documentclass[11pt, letterpaper]{article}
\usepackage[margin=0.8in]{geometry}
\usepackage{booktabs}
\usepackage{array}
\usepackage{tabularx}
\usepackage{hyperref}
\usepackage{parskip}
\usepackage{microtype}
\usepackage{xcolor}
% [cc] T1 + upquote (+ \textquotesingle) so | ' " ` render as literal ASCII symbols.
\usepackage[T1]{fontenc}
\usepackage{upquote}

\newcommand{\nk}{\textcolor{gray}{---}}  % no key (&none)
\newcommand{\gl}[1]{$\to$L#1}            % one-tap switch to layer #1

\title{\textbf{Corne Keymap v22 --- Reference Guide}}
\author{}
\date{2026-11-01}

\begin{document}
\maketitle
\tableofcontents
\newpage

% ------------------------------------------------------------------
\section{Overview}

ZMK keymap for the 42-key Keebmaker Corne (Nice!Nano, split). v22 has
exactly four layers. Layer changes are \emph{persistent one-tap switches}
(no hold-layers), every layer has its own underglow colour, and there are
no transparent keys: a key either has a binding on that layer or does
nothing (shown as \nk).

\begin{center}
\begin{tabular}{clll}
\toprule
\textbf{Layer} & \textbf{Name} & \textbf{Reached by} & \textbf{Underglow} \\
\midrule
L0 & Default (QWERTY) & LT2 or RT2 on L1--L3 & white \\
L1 & Lower  & LT3 on L0 & green \\
L2 & Raise  & LT3 on L1 & blue \\
L3 & Func   & LT3 on L2 & purple \\
\bottomrule
\end{tabular}
\end{center}

% ------------------------------------------------------------------
\section{Notation}

\begin{tabular}{lp{11cm}}
\toprule
\textbf{Name} & \textbf{Meaning} \\
\midrule
LP1--LP3 / RP1--RP3 & Outer pinky column (left / right), numbered top $\to$ bottom \\
LT1--LT3 / RT1--RT3 & Thumb keys, numbered outer $\to$ inner (LT1, RT1 farthest from centre) \\
L0--L3 & Layers by ZMK index \\
M1, M2, M3, M4 & Left Ctrl, Alt/Option, Command, Shift \\
M1b--M4b & Right-hand versions of the same modifiers \\
\gl{n} & One tap switches to layer n (and sets its colour) \\
\nk & No key \\
\bottomrule
\end{tabular}

\begin{center}
\renewcommand{\arraystretch}{1.3}
\begin{tabular}{|c|c|c|c|c|c||c|c|c|c|c|c|}
\hline
LP1 & Q & W & E & R & T & Y & U & I & O & P & RP1 \\ \hline
LP2 & A & S & D & F & G & H & J & K & L & ; & RP2 \\ \hline
LP3 & Z & X & C & V & B & N & M & , & . & / & RP3 \\ \hline
\multicolumn{3}{|c|}{} & LT1 & LT2 & LT3 & RT3 & RT2 & RT1 & \multicolumn{3}{c|}{} \\ \hline
\end{tabular}
\end{center}

% ------------------------------------------------------------------
\section{Layer switching}

\begin{center}
\begin{tabular}{cccc}
\toprule
\textbf{On} & \textbf{LT3} & \textbf{LT2} & \textbf{RT2} \\
\midrule
L0 & \gl{1} & Backspace & Space \\
L1 & \gl{2} & \gl{0} & \gl{0} \\
L2 & \gl{3} & \gl{0} & \gl{0} \\
L3 & \nk    & \gl{0} & \gl{0} \\
\bottomrule
\end{tabular}
\end{center}

LT3 steps forward and stops at L3; either main thumb (LT2/RT2) goes home.
On L1--L3 the first press of LT2/RT2 only returns to L0 (no Backspace /
Space is sent).

% [cc] Usage convention for choosing which thumb to leave L1--L3 with.
\textbf{Which thumb to go home with:} after typing on L1, L2 or L3, if a
space comes next, return with \textbf{RT2} --- it is the Space thumb on L0,
so the next RT2 tap types the space. If no space comes next, return with
\textbf{LT2}. The OLED shows the active layer name (L0--L3).

% ------------------------------------------------------------------
\section{L0 --- Default}

\begin{center}
\renewcommand{\arraystretch}{1.4}
\begin{tabular}{|c|c|c|c|c|c||c|c|c|c|c|c|}
\hline
TAB & Q & W & E & R & T & Y & U & I & O & P & M1 \\ \hline
M2+BSPC & A & S & D & F & G & H & J & K & L & ; & \texttt{\textquotesingle} \\ \hline
M4 & Z & X & C & V & B & N & M & , & . & / & M4b \\ \hline
\multicolumn{3}{|c|}{} & M3 & BSPC & \gl{1} & ESC/M2 & SPC & ENT & \multicolumn{3}{c|}{} \\ \hline
\end{tabular}
\end{center}

\begin{tabular}{lp{11cm}}
\toprule
\textbf{Key} & \textbf{Behavior} \\
\midrule
LP1 TAB & Tab \\
LP2 M2+BSPC & Alt+Backspace (delete previous word on macOS) \\
LP3 / RP3 & Left / right Shift (hold) \\
RP1 M1 & Ctrl (hold) \\
RP2 & \texttt{\textquotesingle} (Shift gives \texttt{"}) \\
LT1 M3 & Command (hold) \\
LT2 / RT2 & Backspace / Space \\
RT3 ESC/M2 & Tap = Escape (sent on release); hold $>$200\,ms = Alt \\
RT1 & Enter \\
\bottomrule
\end{tabular}

% ------------------------------------------------------------------
\section{L1 --- Lower (green)}

\begin{center}
\renewcommand{\arraystretch}{1.4}
\begin{tabular}{|c|c|c|c|c|c||c|c|c|c|c|c|}
\hline
\texttt{`} & 1 & 2 & 3 & 4 & 5 & 6 & 7 & 8 & 9 & 0 & \nk \\ \hline
\nk & SC & SCA & \nk & C-b & \nk & $\leftarrow$ & $\downarrow$ & $\uparrow$ & $\rightarrow$ & \nk & \nk \\ \hline
M4 & M1 & M2 & M3 & \nk & \nk & \nk & \nk & M3b & M2b & M1b & M4b \\ \hline
\multicolumn{3}{|c|}{} & \nk & \gl{0} & \gl{2} & \nk & \gl{0} & \nk & \multicolumn{3}{c|}{} \\ \hline
\end{tabular}
\end{center}

% [cc] Spaces around "+" let LaTeX break these lines (fixes overfull hbox).
SC = macOS screenshot to clipboard (Cmd + Ctrl + Shift + 4); SCA = screenshot
with annotation (Cmd + Shift + 4); C-b = Ctrl + B (tmux prefix). Modifiers on
the bottom row allow e.g.\ Shift+arrows (select), Alt+arrows (word jump),
Cmd+arrows (line start/end), Cmd+1--5.

% ------------------------------------------------------------------
\section{L2 --- Raise (blue)}

\begin{center}
\renewcommand{\arraystretch}{1.4}
\begin{tabular}{|c|c|c|c|c|c||c|c|c|c|c|c|}
\hline
\texttt{\~{}} & \texttt{!} & \texttt{@} & \texttt{\#} & \texttt{\$} & \texttt{\%} & \texttt{\^{}} & \texttt{\&} & \texttt{*} & \texttt{(} & \texttt{)} & \nk \\ \hline
\nk & \nk & \nk & \nk & \nk & \nk & \texttt{[} & \texttt{-} & \texttt{=} & \texttt{]} & \texttt{\textbackslash} & \texttt{\textquotesingle} \\ \hline
M4 & M1 & M2 & M3 & \nk & \nk & \texttt{\{} & \texttt{\_} & \texttt{+} & \texttt{\}} & \texttt{|} & \texttt{"} \\ \hline
\multicolumn{3}{|c|}{} & \nk & \gl{0} & \gl{3} & \nk & \gl{0} & \nk & \multicolumn{3}{c|}{} \\ \hline
\end{tabular}
\end{center}

On the right half each plain symbol has its shifted version directly
below it. Left-hand modifiers allow shortcuts such as Cmd+Shift+[ / ].

% ------------------------------------------------------------------
\section{L3 --- Func (purple)}

\begin{center}
\renewcommand{\arraystretch}{1.4}
\begin{tabular}{|c|c|c|c|c|c||c|c|c|c|c|c|}
\hline
BT\_CLR & BT0 & BT1 & BT2 & BT3 & BT4 & \nk & F1 & F2 & F3 & F4 & \nk \\ \hline
\nk & TOG & EFF & HUI & HUD & BRI & \nk & F5 & F6 & F7 & F8 & \nk \\ \hline
\nk & M1 & M2 & M3 & \nk & BRD & \nk & F9 & F10 & F11 & F12 & \nk \\ \hline
\multicolumn{3}{|c|}{} & \nk & \gl{0} & \nk & \nk & \gl{0} & \nk & \multicolumn{3}{c|}{} \\ \hline
\end{tabular}
\end{center}

\begin{tabular}{lp{11cm}}
\toprule
\textbf{Key} & \textbf{Behavior} \\
\midrule
BT0--BT4 & Select Bluetooth profile 0--4 \\
BT\_CLR & Clear the pairing of the \emph{current} profile (before re-pairing) \\
TOG & Underglow on / off (the only key that turns it on or off) \\
EFF & Next effect (use \emph{solid} so layer colours show cleanly) \\
HUI / HUD & Hue up / down (overridden by the next layer switch) \\
BRI / BRD & Brightness up / down --- persists across layer switches and reboots \\
\bottomrule
\end{tabular}

% ------------------------------------------------------------------
\section{Underglow per layer}

Every layer-switch key is a macro that switches the layer \emph{and}
sets that layer's colour via the custom \texttt{\&rgb\_hs} behavior
(\texttt{src/behaviors/behavior\_rgb\_hue\_sat.c}), which changes hue and
saturation but keeps the current brightness. It never turns the
underglow on or off; TOG controls that. With the underglow off, colours
are still tracked and appear when it is switched on.

\begin{center}
\begin{tabular}{cll}
\toprule
\textbf{Layer} & \textbf{Colour} & \textbf{Hue / Saturation} \\
\midrule
L0 & white  & 0 / 0 \\
L1 & green  & 120 / 100 \\
L2 & blue   & 240 / 100 \\
L3 & purple & 280 / 100 \\
\bottomrule
\end{tabular}
\end{center}

Startup: underglow is off at first boot (\texttt{ON\_START=n}); ZMK
remembers the on/off state and brightness after that.

% ------------------------------------------------------------------
% [cc] v22 diff table added; the v20 vs v21 table is kept below for history.
\section{v21 vs v22 Differences}

\begin{tabular}{p{5.5cm}p{4.5cm}p{5.5cm}}
\toprule
\textbf{Key / feature} & \textbf{v21} & \textbf{v22} \\
\midrule
% [cc] All v22 draft changes were reverted; key behaviour is identical to v21.
All keys & --- & same as v21 \\
\bottomrule
\end{tabular}

% ------------------------------------------------------------------
\section{v20 vs v21 Differences}

\begin{tabular}{p{5.5cm}p{4.5cm}p{5.5cm}}
\toprule
\textbf{Key / feature} & \textbf{v20} & \textbf{v21} \\
\midrule
Return to L0 & LP2 on L1--L3 & LT2 or RT2 on L1--L3 \\
LP2 on L1--L3 & \gl{0} & \nk \\
Underglow & manual only & colour per layer (white/green/blue/purple) \\
Brightness keys & none & BRI / BRD on L3 (G / B) \\
L0 RP2 & M2 (Alt) & \texttt{\textquotesingle} \\
L0 RT3 & ESC & tap ESC / hold M2 (Alt) \\
L2 \texttt{*} & keypad multiply & plain \texttt{*} \\
L2 LP3 & \nk & M4 (Shift) \\
\bottomrule
\end{tabular}

% ------------------------------------------------------------------
\section{Flashing}

\begin{enumerate}
  \item Firmware: \texttt{ysconfig/keyboard/keebmaker\_corne/zmk\_firmware/v22/firmware/}.
  \item Plug in the left half, double-press reset; drive \texttt{NICENANO} appears.
  \item Copy \texttt{corne\_left-nice\_nano\_\_zmk-zmk.uf2} onto it; it reboots.
  \item Repeat for the right half with the right \texttt{.uf2}.
\end{enumerate}

\end{document}
